\scanpage{186}\mspage{72}
Since
\handwritten{\[
O_V(-E_{\varphi})[nE]\cong O_V[(n-1)E]
\]}
then, having taken $n''\geq n'$, we obtain from (14) that, at $n\geq n''$
\handwritten{\[
\begin{aligned}
&\dim H^1(V,O_V[nE])=\\
&=\dim H^1(V,O_V(-\mathscr{E}_{\varphi})[nE])=\dim H^1(V,O_V[(n-1)E]).
\end{aligned}
\tag*{(15)}
\]}

\margin{p. 208}It is evident that $\unsure{\varphi|\omega}^{n}=\{\varphi_i^{n}\}$ is a cross-section of the group $H^0(V,O_V[nE])$. Therefore, the common zeros of the elements of the group $H^0(V,O_V[nE])$ are contained among the zeros (of) $\varphi^n$, that is, all of them lie on the subvariety $\mathscr{E}_{\varphi}$.

Let us show that, at $n\geq n_{\mathscr{E}}$ there is no common zeros of the cross-sections of $H^0(V,O_V[nE])$ on $\mathscr{E}_{\varphi}$. Consider an exact sequence
\handwritten{\[
0\to O_V(-\mathscr{E}_{\varphi})\to O_V\xrightarrow{r}O_{\mathscr{E}_{\varphi}}\to 0
\tag*{(16)}
\]}

The homomorphism $r$ is described in the following way: if $P\notin\mathscr{E}_{\varphi}$, then $r$ in the stalk over this point is a zero homomorphism; if $P\in\mathscr{E}_{\varphi}$, then $r$ over $P$ is
% archive sheet number: 185
